\section{The coefficient offset $\Delta$}\label{sec:delta}

\subsection{Definition}
Let $\Delta$ denote the amount by which the realized kinetic coefficient exceeds the
requested one. From Eq.~\eqref{eq:tau} with $C_{\rm corr}=0$,
\begin{equation}
\Delta \;\equiv\; \bbare-\bsub \;=\; a_1a_2\,\epsilon
\left(\frac{1}{\texttt{diff\_sub}}+\frac{1}{\texttt{dif\_vap}}\right)
\frac{\rho_{\rm ice}}{\rho_{vs}(T)} ,
\label{eq:delta}
\end{equation}
expressed in the unscaled units of $\bsub$ (s\,m$^{-1}$). The bracket is the pair of $a_2$
terms carried by $\tau_{\rm sub}$; the factor $a_1$ converts the bracket into a kinetic
coefficient, and $\rho_{\rm ice}/\rho_{vs}$ removes the scaling used internally by the
solver. For the wedge configuration ($\epsilon=8.584\times10^{-7}$\,m,
$T=-20\,^\circ$C):
\begin{center}\small\begin{tabular}{lrr}\toprule
 & value [s\,m$^{-1}$] & share of $\Delta$\\ \midrule
thermal, $a_1a_2\epsilon\,\rho_{\rm ice}/(\rho_{vs}\texttt{diff\_sub})$ & $9.44\times10^{4}$ & 73.7\,\%\\
vapor, \ \,$a_1a_2\epsilon\,\rho_{\rm ice}/(\rho_{vs}\texttt{dif\_vap})$ & $3.37\times10^{4}$ & 26.3\,\%\\ \midrule
$\Delta$ & $1.282\times10^{5}$ & \\
$\bsub$ (requested) & $5.9216\times10^{5}$ & \\
$\bbare=\bsub+\Delta$ (realized) & $7.2036\times10^{5}$ & \\ \bottomrule
\end{tabular}\end{center}

\subsection{Consequences}
Three properties of Eq.~\eqref{eq:delta} determine its practical significance.

\emph{(i) The offset is additive.} The resulting shortfall is $\bsub/(\bsub+\Delta)$ and
therefore depends on $\bsub$ itself; it is not a fixed multiplicative factor. The
parameter sweep of \S\ref{sec:ladder} exhibits this directly, the ratio varying from
$0.32$ to $0.95$ across four values of $\bsub$. Consequently, any correction of the form
``rescale $\beta$ by the observed ratio'' is calibrated at a single parameter point and is
inaccurate elsewhere (\S\ref{sec:fix}).

\emph{(ii) The offset scales as $\epsilon$ and as $1/\rho_{vs}(T)$.} Coarser interfaces and
lower temperatures both increase it, and neither quantity is freely adjustable: $\epsilon$
is fixed by resolution requirements and $T$ by the problem under study. For the Molaro
configuration ($\epsilon=1.2\times10^{-7}$\,m, $\alpha_c=0.1$), $\Delta/\bsub=0.23$; for
the wedge configuration, $0.22$; at $-60\,^\circ$C, $20.7$.

\emph{(iii) The offset imposes a lower bound on the realizable kinetic coefficient.} Since
$\beta_{\rm realized}=\bsub+\Delta\ge\Delta$ for any positive $\bsub$, and
$\beta\propto1/\alpha_c$, there exists an upper bound on the attachment coefficient that no
choice of input can exceed:
\begin{center}\small\begin{tabular}{rrrr}\toprule
$T$ [$^\circ$C] & $\rho_{vs}$ [kg\,m$^{-3}$] & $\Delta$ [s\,m$^{-1}$] & max.\ attainable $\alpha_c$\\ \midrule
$-60$ & $8.885\times10^{-6}$ & $1.225\times10^{7}$ & 0.0006\\
$-40$ & $1.055\times10^{-4}$ & $1.031\times10^{6}$ & 0.0077\\
$-20$ & $8.487\times10^{-4}$ & $1.282\times10^{5}$ & 0.0619\\
$  0$ & $5.047\times10^{-3}$ & $2.156\times10^{4}$ & 0.3684\\ \bottomrule
\end{tabular}\end{center}
(evaluated at $\epsilon=8.584\times10^{-7}$\,m; halving $\epsilon$ halves $\Delta$ and
doubles each bound). The Molaro configuration is affected directly: a requested
$\alpha_c=0.1$ is realized as $0.081$, and in the limit $\alpha_c\to1$, appropriate to an
attachment-limited regime, $\Delta$ exceeds $\bsub$ by a factor of $2.3$.

\section{Vanishing of the correction for a one-sided diffusion field}\label{sec:inner}

\subsection{Nature of the correction}
The thin-interface correction is not a physical effect. At finite $\epsilon$ the
diffuse-interface formulation introduces a spurious contribution to the sharp-interface
kinetic coefficient, arising because $\sigma$ varies across the interfacial region in
which the phase-change source is distributed. $\tau_{\rm sub}$ is inflated by the
anticipated magnitude of that contribution so that the two cancel and the model realizes
$\bsub$. The construction is standard; the question addressed here is whether the present
model in fact generates the contribution for which it is being compensated.

\subsection{The solvability condition}
Introduce the signed normal distance $\zeta$ (positive into the vapor) and the stretched
inner coordinate $\eta=\zeta/\epsilon$. In the inner region the ice phase field, written
$\phi$ here and $\phi_i$ elsewhere in the project, is expanded as
\begin{equation}
\phi=\phi_0(\eta)+\epsilon\,\phi_1(\eta,s,t)+O(\epsilon^2),
\end{equation}
in which the subscript denotes order in $\epsilon$ and is not a phase index. The leading
term satisfies $\partial^2_{\eta\eta}\phi_0-f'(\phi_0)=0$ and is therefore the planar,
stationary equilibrium profile $\phi_0=\tfrac12\left(1+\tanh(\eta/2)\right)$; curvature,
interface motion and the driving field enter at $O(\epsilon)$ through $\phi_1$.

At $O(\epsilon)$ the linearized operator is singular, and the Fredholm alternative,
obtained by projection onto the translation mode $\partial_\eta\phi_0$, yields the
interface condition. Writing $\sigma=\sigma_{\rm int}+\delta\sigma(\eta)$, where
$\sigma_{\rm int}$ is the outer solution extrapolated to $\zeta=0$, the coefficients follow
from
\begin{equation}
\int(\partial_\eta\phi_0)^2\,\mathrm{d}\eta=\tfrac16,\qquad
\int g(\phi_0)\,\partial_\eta\phi_0\,\mathrm{d}\eta=\tfrac1{30},\qquad
a_1=\frac{1/6}{1/30}=5 ,
\end{equation}
evaluated numerically to six significant figures and equal to the value of $a_1$ used by
the solver. The correction depends on a third integral,
\begin{equation}
I=\int g(\phi_0)\,\partial_\eta\phi_0\;\delta\sigma(\eta)\,\mathrm{d}\eta ,
\label{eq:Iint}
\end{equation}
and vanishes if and only if $I$ does.

The weight $g(\phi_0)\,\partial_\eta\phi_0$ is even in $\eta$, since
$\phi_0(-\eta)=1-\phi_0(\eta)$ renders $g=\phi_0^2(1-\phi_0)^2$ even and
$\partial_\eta\phi_0=\tfrac14\,\mathrm{sech}^2(\eta/2)$ even. The integral $I$ therefore
samples only the even part of $\delta\sigma$.

\subsection{Symmetric diffusivity}
For $D$ equal in both phases and the flux partitioned as $\pm j/2$, the quasi-static inner
balance $\partial^2_{\eta\eta}\rho_v\propto\partial_\eta\phi_0$ integrates to
$\partial_\eta\rho_v\propto(\phi_0-\tfrac12)$, which is odd. Hence
$\delta\sigma\propto\Psi(\eta)=\ln\cosh(\eta/2)-|\eta|/2+\ln2$ is even and decays to zero
in both limits. An even deviation against an even weight gives $I\neq0$; this is the
origin of the term $-a_1a_2\epsilon/D$ of Karma \& Rappel (1998).

\subsection{One-sided diffusivity}
The vapor equation of the present model carries the diffusivity $D_v\phi_a$, which vanishes
in the ice. Neglecting storage, which is justified by
$v_n\epsilon/(\xi_v D_v)=2.4\times10^{-11}$, the inner balance is
\begin{equation}
\partial_\eta\!\left(D_v\phi_a\,\partial_\eta\rho_v\right)
=-\epsilon\,v_n\rho_{\rm ice}\,\partial_\eta\phi_0 .
\end{equation}
A single integration, with the constant fixed by the absence of flux into the ice at
$C=\epsilon v_n\rho_{\rm ice}$, gives
\begin{equation}
D_v\phi_a\,\partial_\eta\rho_v=\epsilon v_n\rho_{\rm ice}(1-\phi_0)
=\epsilon v_n\rho_{\rm ice}\,\phi_a
\;\;\Longrightarrow\;\;
\partial_\eta\rho_v=\frac{\epsilon v_n\rho_{\rm ice}}{D_v}
\quad\text{wherever }\phi_a\neq0 .
\end{equation}
The factor $\phi_a$ cancels, and the resulting constant is precisely the outer Stefan
gradient, since $D_v\partial_\zeta\rho_v|_{\rm int}=v_n\rho_{\rm ice}$ yields the same
value. The inner solution is thus the continuation of the outer solution through the
interfacial region, so that
\begin{equation}
\delta\sigma\equiv0
\;\;\Longrightarrow\;\; I=0
\;\;\Longrightarrow\;\; C_{\rm corr}=0 .
\end{equation}
The integrand vanishes identically rather than by cancellation of contributions: since
$\delta\sigma$ is null rather than merely odd, $I=0$ for any weight, and no evaluation of
the prefactor is required. Deep within the ice the floor \code{air\_lim}$=10^{-6}$ on
$\phi_a$ breaks the cancellation, but $g(\phi_0)$ there is of order $10^{-12}$ and the
contribution is negligible.

\subsection{Limitations of the present derivation}
The vanishing result above is insensitive to prefactors and is correspondingly robust. The
symmetric branch is less well established here: constructing $\delta\sigma=\Psi$ as above
and evaluating Eq.~\eqref{eq:Iint} gives $a_2=-6I=0.087$, against the accepted value
$0.158$. The present treatment therefore omits $O(\epsilon)$ contributions, plausibly the
curvature term, the back-reaction of $\phi_1$, or the matching of the flux partition. It is
not offered as a derivation of $a_2$. What is established is the structural contrast
between an even, non-vanishing deviation and an identically null one; the magnitude is
supplied by measurement.

\subsection{The thermal channel}
The thermal term is not affected by the one-sidedness argument: $k$ is interpolated and
non-zero in both phases, so a correction does exist for that channel. It is nonetheless
misspecified, because \code{a2*eps/diff\_sub} presumes that latent heat drives $\sigma$ at
unit strength. Temperature influences $\sigma$ only through $\rho_{vs}(T)$, so the channel
carries a Clausius--Clapeyron factor. Evaluating
$\mathrm{d}\sigma/\mathrm{d}(v_n\epsilon)$ for each channel:
\begin{center}\small\begin{tabular}{llr}\toprule
channel & coefficient & value\\ \midrule
vapor & $1/D_v$ & $5.27\times10^{4}$\\
      & as coded: $1/\texttt{dif\_vap}$ & $4.59\times10^{4}$ \ (correct form; $I=0$)\\[2pt]
thermal & $(\mathrm{d}\rho_{vs}/\mathrm{d}T)\,L_{\rm sub}/k$ & $1.0\times10^{2}$ ($k_{\rm ice}$) -- $1.2\times10^{4}$ ($k_{\rm air}$)\\
        & as coded: $1/\texttt{diff\_sub}$ & $1.29\times10^{5}$ \ (over-weighted by 11--1300)\\ \bottomrule
\end{tabular}\end{center}
Since \code{ThermalCond} interpolates between the two conductivities, any admissible choice
renders this term negligible. This misspecification is independent of the preceding
argument.

\subsection{Scope of the measurement}
\begin{center}\small\begin{tabular}{lll}\toprule
spurious term & form & constrained by the sweep?\\ \midrule
kinetic correction & $\propto v_n$ & yes; $\le0.1\,\%$ of $\bbare$\\
solute trapping & $\propto\epsilon v_n$ (shifts $\sigma_{\rm int}$) & yes; also $\propto v_n$, absorbed into fitted $\beta$\\
surface diffusion & tangential flux along the interface & no; distinct operator\\ \bottomrule
\end{tabular}\end{center}
The first two are excluded. The third is not observable in a centerline $v_n$ measurement
and is not constrained by any result reported here. The surface excess of the diffusivity
weight, $\int[(1-\phi_0)-\Theta(-\eta)]\,\mathrm{d}\eta$, vanishes exactly by the same
antisymmetry, which suggests that the spurious surface diffusion is likewise suppressed;
this is an indication rather than a calculation. The point is material because the model
contains no surface-diffusion term by construction, so a spurious one would not be evident
from the source, and the Molaro analysis separately invokes physical surface diffusion.
The two must be distinguished.
